\subsection{Refined Equilibrium Characterization}

We make the cutoff structure explicit in terms of posterior beliefs. Let the posterior mean be:
\begin{equation}
\hat{\theta}_i = w_0 \mu + w_x x_i + w_s s.
\end{equation}

Since $\hat{\theta}_i$ is monotone in $x_i$, the cutoff strategy in signal space maps into a cutoff in posterior space. There exists $\hat{\theta}^*$ such that agents attack iff:
\begin{equation}
\hat{\theta}_i \geq \hat{\theta}^*.
\end{equation}

This implies a corresponding cutoff in signals:
\begin{equation}
x^* = \frac{\hat{\theta}^* - w_0 \mu - w_s s}{w_x}.
\end{equation}

\subsection{Sharpened Comparative Statics}

\textbf{Proposition 5 (Refined).} An increase in public signal precision $\tau_s$ lowers the posterior cutoff $\hat{\theta}^*$.

\textit{Proof.}

The indifference condition is:
\begin{equation}
\pi(\hat{\theta}^*) \cdot B = C.
\end{equation}

Higher $\tau_s$ increases the weight on the common signal and reduces posterior variance. This steepens the mapping from fundamentals to aggregate behavior. As a result, for a given benefit-cost ratio, the required posterior threshold for coordination decreases. Formally, $\frac{\partial \hat{\theta}^*}{\partial \tau_s} < 0$. \hfill $\square$

\textbf{Proposition 6 (Refined).} Public signal precision increases the sensitivity of aggregate attacks to fundamentals.

\textit{Proof.}

Aggregate attack is:
\begin{equation}
R(\theta) = 1 - \Phi\left(\frac{x^* - \theta}{\sigma_x}\right).
\end{equation}

Using the chain rule:
\begin{equation}
\frac{dR}{d\theta} = \phi\left(\frac{x^* - \theta}{\sigma_x}\right) \cdot \frac{1}{\sigma_x}.
\end{equation}

Since $x^*$ decreases in $\tau_s$, more mass lies near the threshold region, increasing local sensitivity. \hfill $\square$

\subsection{Global Games Mapping}

The model maps to a standard global games environment with endogenous public precision. The key sufficient statistic is total precision:
\begin{equation}
\tau = \tau_x + \tau_s(q_s).
\end{equation}

As $\tau_s$ increases, the dispersion of posterior beliefs shrinks:
\begin{equation}
\text{Var}(\hat{\theta}_i) = \frac{1}{\tau_x + \tau_s + \tau_\theta}.
\end{equation}

Lower dispersion implies tighter clustering of actions, strengthening coordination.

\subsection{Welfare Threshold (Refined)}

\textbf{Proposition 8 (Refined).} There exists a unique threshold $\theta^*$ such that:
\begin{equation}
\frac{dW}{dq_s} > 0 \quad \text{for } \theta < \theta^*, \quad \frac{dW}{dq_s} < 0 \text{ for } \theta > \theta^*.
\end{equation}

\textit{Proof.}

Differentiate welfare:
\begin{equation}
\frac{dW}{dq_s} = \frac{\partial W^{eff}}{\partial q_s} - \frac{\partial W^{frag}}{\partial q_s}.
\end{equation}

The efficiency term is decreasing in transaction costs and always positive. The fragility term is increasing in crisis probability, which is increasing in $\tau_s(q_s)$ and thus $q_s$. At low $\theta$, $\frac{\partial W^{frag}}{\partial q_s}$ is small; at high $\theta$, it dominates. Continuity ensures a unique crossing. \hfill $\square$

\subsection{Interpretation}

The refined results make clear that the core mechanism operates through the compression of beliefs. Stablecoins reduce disagreement, which is beneficial for price discovery but destabilizing in environments with strategic complementarities.
